%% kapa: exam=1 round=31 year=2020 period=1 form=A subject=civil_law subject_ko=민법 seq=26 number=26
%% source: Q-Net 감정평가사 기출문제 (cst00309) artlSeq=5207219 file=extracted/1008/5207219/A형/A형/2020년도 제31회 감정평가사 1차 시험_1교시_A형.pdf
%% conversion: kordoc 4.24.51 -> tools/gen1.py (auto); review=proofread-sonnet; flags=none
\begin{question}{26}
선의취득에 관한 설명으로 옳지 않은 것은? (다툼이 있으면 판례에 따름)
\begin{choices}
  \item 점유권과 유치권은 선의취득할 수 없다.
  \item 점유개정의 방법으로 양도담보를 설정한 동산소유자가 다시 제3자와 양도담보설정계약을 맺고 그 동산을 점유개정으로 인도한 경우, 제3자는 양도담보권을 선의취득하지 못한다.
  \item 인도가 물권적 합의보다 먼저 이루어진 경우, 선의․무과실의 판단은 인도를 기준으로 한다.
  \item 선의취득자는 임의로 소유권취득을 거부하지 못한다.
  \item 선의취득자는 권리를 잃은 전(前)소유자에게 부당이득을 반환할 의무가 없다.
\end{choices}
\end{question}
\answer{3}
